\subsection{The numerical regions and the induction}
\label{j:finite:assembly}
The preceding lemmas reduce each finite index to a closed cover of
moment and probability parameters. Table~\ref{j:finite:region-table}
shows how the methods meet. All displayed intervals are closed; adjoining
rows overlap at their endpoints. The target of each Fourier region is
$T=4097/10000<\CE$. The affine bound closes the indicated upper tail;
at $\rho=7361/4000$ it is already at most $819/2000<T$.
\begin{table}[htbp]
\centering\small
\caption{The complete third-moment cover for indices five through thirty-nine.
The count column includes necessary-condition exclusions and, where specified,
the signed corner or finite-sine comparison.}
\label{j:finite:region-table}
\begin{tabular}{@{}cP{31mm}P{66mm}c@{}}
\toprule
$n$&Count region&Fourier complement and required support&Affine tail\\\midrule
5&$[1,6/5]$, with the positive corner below
 &$[6/5,7361/4000]$; the six support rows of
 Proposition~\ref{j:finite:supportrows}&$\rho\ge7361/4000$\\
6--9&$[1,237/200]$
 &$[237/200,6/5]$, radius $13/5$; then the five support rows
 through $7361/4000$&$\rho\ge7361/4000$\\
10--14&$[1,23/20]$
 &$[23/20,71/50]$, radius $3$; $[71/50,15/8]$, unrestricted
 &$\rho\ge7361/4000$\\
15--19&$[1,23/20]$
 &$[23/20,33/25]$, radius $3$; $[33/25,19/10]$, unrestricted
 &$\rho\ge19/10$\\
20--24&$[1,9/8]$
 &$[9/8,6/5]$, radius $3$; $[6/5,19/10]$, unrestricted
 &$\rho\ge19/10$\\
25--34&$[1,6/5]$
 &$[6/5,19/10]$, unrestricted&$\rho\ge19/10$\\
35--39&$[1,6/5]$
 &$[6/5,19/10]$, unrestricted&$\rho\ge19/10$\\\bottomrule
\end{tabular}
\end{table}

The count regions split at $11/10$; at indices five through nine they
also split at $23/20$. At five the exact signed cover is as follows.
Put $b_2=84729828287999/(250000000000000)$ and
$b_3=35139006743001/10^{14}$. The signed-moment inequality gives
$|\beta|\le b_2$ up to $\rho=237/200$ and $|\beta|\le b_3$ up to
$\rho=6/5$: at each endpoint the rational test is
$(1+b_i^2)^2\ge\rho_i^2(1-b_i^2)$, with the difference increasing
in $b_i\ge0$. The count cover contains every feasible signed moment
for $1\le\rho\le477/400$, and every such moment with
$\beta\le17/50$ for $477/400\le\rho\le6/5$.
The remaining closed rectangle is exactly
Proposition~\ref{j:finite:gapfree5-result}; its positive sign is retained.
The splits at $237/200$, $467/400$, $477/400$, $1/5$ and $3/10$
are refinements of this union, not additional assumptions.

\paragraph{What the numerical assertion checks.}
The finite acquisition theorem~\ref{j:local:acquisition} transports
one actual marginal through every update. Support is proved before
higher moments are used, and an initial raw remainder before defect
feedback. Its simultaneous state includes the score, support, moments,
defect and signed threshold. A failed analytic gate gives no exclusion.
Each state must instead reach a necessary-condition contradiction, a
complete Fourier bound, or a complete count forest.

Write $P_j=[j/100,(j+1)/100]$, $0\le j<100$.
Every probability panel is closed. When several moment cells export
to it, upper bounds are maximized and lower bounds minimized over
\emph{all} those cells. Thus seams retain the source law. A count
forest covers every original root $k=0,\ldots,n$, including all signed
phase intervals. Missing panels are justified exclusions; refinements
must cover their entire parent. The exact partitions, score/HIGH caps,
refinement schedules and forest counts are certificate data. Their
logical content is the complete-domain assertion below, not the number
of successful files.

The following application-specific prerequisites are retained.
At five the first own-score bound is support-free; later Fourier inputs
use Proposition~\ref{j:finite:supportrows}, including every narrower
signed profile. For six through nine, the wider Fourier domain
$[23/20,6/5]$ uses radius $13/5$, while its low-band resource uses
radius $5/2$. At fifteen through nineteen, support radius three on
$\rho\le33/25$ precedes the own-score and HIGH bounds; at twenty,
radius three on $\rho\le6/5$ precedes them. At twenty-one through
twenty-four the own-score bounds are unrestricted. The twenty through
twenty-four initializers require no earlier numerical seed, and each
reused coarse count contract at twenty is admitted only after full
simultaneous containment has been established. A cellwise Fourier
bound evaluates the complete compensated expression, not a presumed
scalar HIGH ceiling. Each unrestricted region includes the complete
auxiliary-moment domain.

At $26$--$30$ and $32$--$34$, respectively, the acquired floors are
$N=25$ and $N=31$. Transfer retains $nD\le Nd$ as well as the raw
remainder bound. All binomial probabilities, normal arguments and
support tests still use the actual integer $n$. Local refinements of
a parent replace only the part they cover. At $26$--$30$ the surviving
second-band panels $P_{32},\ldots,P_{42}$ are closed by the finite-sine
inequality; the unrestricted Fourier complement makes a third count
band at twenty-five unnecessary. All omitted count panels, including
the two sign-mass exclusions at twenty and the weighted-budget
exclusions at $26$--$30$, have their own necessary-condition checks.

Here the five weighted exclusions have a direct explanation.
For $(n,j)=(30,52)$ and for $27\le n\le30$, $j=53$, the acquired
pairs $(d,\Theta)$ in $D\le25d/n$, $\theta\le\Theta$ are respectively
$(3/10,1/500)$ and $(5/16,1/2500)$. Put
$s=316227766/10^8<\sqrt{10}$ and $l=j/100$. Since
$p\ge l>p_*$, $\sqrt{p(1-p)}\le1/2$ and
$(1-\theta)^{3/2}\ge(1-\Theta)^2$, the joint budget gives
\[
 D\ge(12+4s)(1-\Theta)^2(l-(4-s)/2)^2
 >\begin{cases}251/1000,&j=52,\\38/125,&j=53.\end{cases}
\]
These contradict $D\le1/4$ at $(30,52)$ and
$D\le125/432<38/125$ at $j=53$, $n\ge27$.
The separate exclusions of $P_{36}$ and $P_{34}$ at twenty use
the sign-mass inequality on the respective bands, including endpoints.

\paragraph{The prerequisite at thirty-nine.}
Sharpness through 38 and the recurrence give
\[
 39(C_{39}-\CE)<55887/190000<3/10,\qquad C_{39}<.4226,
\]
with both deleted scores below $.4253$. Before the sharp-39 acquisition,
apply the separate comparison with actual $n=39$, $K=741/4000$,
predecessor allowance $.33$ and excess allowance $1/2$.
Under $C_{39}\ge\CE+K/39$, its support input is
$(N,L,K,q,R)=(39,39,741/4000,.33,3.84)$, and the positive charge
$K\rho_l$ remains in the resource. Its full acquisition excludes the
probability complement of $[.302,.458]$.
Nearest count covers all phases for $0<\lambda\le.08$; on
$[.08,.10]$ use nearest count for $|u|\ge.19$ and whole count for
$|u|\le.19$; on $[.10,.12]$ use nearest count for $|u|\ge.27$
and whole count for $|u|\le.27$;
above $.12$ whole count covers every phase. Each band is clipped to its
own noise ceiling and includes its seams. Zero noise is Bernoulli.
The signed allowance tests of Section~\ref{j:local:allowance} exclude
the hypothesized score. Hence $C_{39}<.41449$, since
$\CE<.40974$ and $K/39=.00475$.

At 35--38, the entry bound
$C_n<41/100+4177/(13824n)$ supplies the own-score caps; at 39 use
the bound just proved. The initial support radii are $39/10$ at 35
and $19/5$ at 36--39. The independent support bound $96/25<39/10$
licenses the first, and all five initial and resource calculations
use the actual index. Any local support refinement is proved on its
whole source interval before export. Thus sharpness at 39 is never
used to establish its own entry cap.

\begin{proposition}[Complete numerical finite base through thirty-nine]
\label{j:finite:numerical539}
The support, Fourier, initialization, resource and count inequalities
in the preceding closed cover hold on their complete domains, including
the separate positive-39 prerequisite. Every exported row implies its
entire original count contract. The necessary exclusions and retained
forests cover every original moment and probability panel.
\end{proposition}
This assertion combines acquisition for one marginal, directed
containment of every scalar and signed-phase condition, count acceptance
at the actual integer, and the complete closed union. In particular,
if $q_i(H)$ are the acquired phase polynomials and
$\widetilde q_j(H)$ the contract polynomials, the join proves, for each
$i$, some $j$ with $q_i\le\widetilde q_j$ throughout that signed interval.
Endpoints and the possible interior minimum of the quadratic difference
suffice. Hence $\max_iq_i\le\max_j\widetilde q_j$. Selecting one row
retains the whole contract's actual index and moment domain.

\begin{proof}[Derivation of $C_n\le\CE$ for $5\le n\le39$]
Choose the least failing index and an attained positive closed maximizer.
Its predecessor is sharp, so contact loss is zero and every larger
retained allowance is valid. The table covers every $\rho\ge1$.
Fourier and signed-corner regions have score below $T<\CE$.
Otherwise acquisition retains the same law in an exported panel.
It cannot satisfy a necessary-condition exclusion; the directed join
therefore puts it in a complete count root, whose bound is below $\CE$.
At 39 the positive comparison is applied first. Every case contradicts
the assumed failure.
\end{proof}

\subsection{The actual indices forty, forty-one and forty-two}
\label{j:finite:4042}
A first failure here has sharp predecessors through two deletions.
The recurrence gives
\[
 0<n(C_n-\CE)<1433/4875<3/10,\qquad C_n<167/400<.4226.
\]
Thus the deleted bounds are below $.4253$, and the weaker allowances
$n(C_n-\CE)\le.54$ and contact loss $.19$ are valid.
The floor-40 argument preserves $\mathcal R_n\le r$ and $nD\le40d$;
its support test uses both endpoints 40 and 42 and proves radius $96/25$
before higher moments are used. Complete acquisition leaves left rows
covering $[.282,.38]$ and central rows covering $[.38,.476]$; their
complements are excluded. Confinement starts with widths
$.55,.575,.6$ on the left and $.55$ centrally, and cubic constants
30 and 24, with every width and positive-slope gate checked.

The central phase cover uses nearest count for $0<\lambda\le.08$,
then nearest count for $|u|\ge.19$ and whole count for $|u|\le.19$
on $[.08,.10]$, and nearest count for $|u|\ge.27$ and whole
count for $|u|\le.27$ on $[.10,.12]$;
whole count applies above $.12$.
The left cover uses nearest count through $.0925$, then nearest
count for $|u|\ge.16$ and whole count for $|u|\le.16$ on
$[.0925,.10]$, and nearest count for $|u|\ge.25$ and whole count
for $|u|\le.25$ on $[.10,.12]$; then nearest count for $u\le-.27$
and whole count for $u\ge-.27$ on $[.12,.135]$; above $.135$
use whole count. Each closed band is clipped to its acquired noise
ceiling. Zero noise is Bernoulli.

These three indices admit a positive finite-binomial reserve. For
$[a,b]\subset(0,1/2)$ put
$F_j(p)=\Pr(\operatorname{Bin}(n,p)\le j)$ and
$z_j(p)=(j-np)/\sqrt{np(1-p)}$. Both decrease with $p$, while
$\sqrt{np(1-p)}/(1-2p(1-p))$ increases. Therefore both sides of
every jump are bounded by
\begin{equation}\label{j:finite:binomialslack}
 \frac{\sqrt{nb(1-b)}}{1-2b(1-b)}
 \max\{F_j(a)-\Phi(z_j(b)),\Phi(z_j(a))-F_{j-1}(b),0\},
 \qquad F_{-1}=0.
\end{equation}
Between jumps the CDF is constant, so $j=0,\ldots,n$ suffices.
Exact endpoint probabilities and outward Gaussian bounds on a complete
closed subdivision of $[.28,.48]$ give a positive uniform lower bound
$\epsilon_p$ for $\CE-C^{\rm bin}_{n,p}$ on each parent, at the actual
$n=40,41,42$.

Retain both signs in the finite corrections:
$I/n\le\max(I/40,I/42)$ and
$\sqrt n\pi_kJ\le\max(b_0J,J)$ when $b_0\le\sqrt n\pi_k\le1$.
If $S$ bounds $\rho_B(C-C_B)/\lambda$ on a positive-noise band
$[a,b]$, it suffices that $S<\rho_{B,l}\epsilon_p/b$.
For $S<0$ this is immediate; otherwise use $\lambda\le b$ and
$\rho_B\ge\rho_{B,l}$. The scalar bounds are convex in positive
noise, being affine terms plus nonnegative reciprocals of positive
affine terms; their two endpoint tests cover the band. No assertion
for all $n\ge40$ follows from this finite calculation.

\begin{proposition}[Numerical floor-40 assertion]\label{j:finite:numerical4042}
The preceding acquisition, two-endpoint support, actual-index binomial
reserve and complete signed local covers satisfy their strict tests
at $n=40,41,42$. Each local row consumes its acquired simultaneous
profile. Consequently $C_{40},C_{41},C_{42}\le\CE$.
\end{proposition}
\begin{proof}
The base through 39 and a least failure supply all entry premises.
Acquisition and confinement preserve its witness. The local comparison
with the signed finite corrections excludes every positive-noise case;
the Bernoulli theorem excludes zero noise.
\end{proof}

\begin{theorem}[The finite base]\label{j:finite:base}
For every integer $1\le n\le42$ and every standardized iid marginal
with finite third absolute moment, $\sqrt n\,\Delta_n(P)\le\CE\rho(P)$.
\end{theorem}
\begin{proof}
The starting bounds and the four-summand comparison establish indices
one through four. Proposition~\ref{j:finite:numerical539} and its
same-law induction establish five through thirty-nine;
Proposition~\ref{j:finite:numerical4042} supplies the last three indices.
\end{proof}
